% LTeX: language=en-GB enabled=true

\documentclass[../main.tex]{subfiles}

\begin{document}
\newenvironment{qanda}{\setlength{\parindent}{0pt}}{\bigskip}
\newcommand{\Q}{\bigskip\bfseries Q: }
\newcommand{\A}{\par\textbf{A:} \normalfont}

\subsection{Introduction}

A backup utility is an essential software program for any computer system, allowing users to keep their files and data safe in the case of hardware failure, or various possible software issues. It is particularly important that these files are stored externally, on another device, in case of such hardware failure. Therefore, I will create a backup utility for storing files on an external server from a client PC. This will require various components, such as encrypting the files stored on the server, encrypting the transmission of files between client and server, and compressing the files being transferred. 

The main benefits of the program I will create over existing systems are that my utility will be fully open source, so the function of the program can be independently inspected by any user, and that it will be tailored to my client’s exact needs: many existing applications are either closed source, or too bloated with unnecessary features and unintuitive interfaces. Instead, this will be a simple, easy-to-use command line tool.

I will research this through two main methods: a discussion with my client, and using reliable external sources. The latter may come from reputable books, websites, papers, or journals, and will be cited. 

\subsection{Client Overview}

My client will be a classmate – Sam Yarasani – who informed me he is looking for such software. Therefore, I will tailor this program for his exact needs, and I will establish his specific requirements through a client interview. This has the benefit over existing systems of being in exactly the format he would like, without additional features which add complexity and bloat to the final program. Plus, it will have the exact user interface he would like.

\subsection{Client Interview} \label{ClientInterview}

Below is a record of the transcript of my first interview with my client. I will use this to inform the design requirements for this backup utility. 

\begin{qanda}

\Q I understand that you are looking for a backup utility? 
\A Yes, I'd like to be able to create external backups of my system on a separate device. 

\Q What devices would you like to create backups of? 
\A My laptop and desktop PC. 

\Q What operating system are those devices running? 
\A Arch Linux, and Ubuntu 22.04 respectively, but I would like this to work on any modern Linux distribution - I may want to change in the future! 

\Q Where would you like the backups to be stored? 
\A I have a fairly old laptop which I use to run a few home server utilities. I’d like to also store the backups here, if possible. 

\Q Can you be more specific about the laptop? 
\A It's around 15-20 years old, I'm not exactly sure, with 4GB of RAM. Also, it's running Ubuntu server. 

\Q Would you like a third-party server in a data centre to store the backups instead? 
\A Absolutely not. The key here is that this is all fully local to me, and that I am the only one with access to the backups. Many current backup utilities lack this! 

\Q Would you want the files stored to be encrypted in transmission, and while stored? 
\A It is crucial for the files to be hidden while being transferred, as I do not want any data being intercepted by others. However, I don't need the files to be encrypted while being stored on the server, as I own that server and know that it is kept safe.

\Q Do you have any other priorities? 
\A File transfers should be as quick as possible, and as little space as possible should be used on the server device. Oh, and I want to be sure that transfers have been successful. 

\Q What type of files will you be backing up? 
\A Well, I will be using this for my whole system, so all sorts of files. I guess there will be text files, like system configuration files, image files, video files, and raw binary executables. 

\Q Okay. The non-text files you were talking about, are they already compressed? 
\A Yes actually, almost all of my audio and video files are already compressed. 

\Q Very good, thanks. What about the user interface, what form should this take? 
\A I want a very simple command line interface, with no messy graphical elements. It could have an interactive mode, but there should certainly be a way to use it non-interactively, with just command line flags, so that I can use it in scripts. 

\Q Finally, is it important to you for this project to be finished by a certain date? 
\A Absolutely. It is essential for this to be finished by the end of March 2026. This is one of my key priorities, since I will start work on an important project in April, and I need to ensure I have safe backups of my files by its start.

\end{qanda}

I have summarised my key findings from this interview below:
\begin{itemize}
    \item The program should work on Linux systems,
    \item The program should be able to comfortably run on my client's devices, without putting much strain on them,
    \item I must also code the server storing the files, which should run on a device my client also owns,
    \item The files must be encrypted in transmission,
    \item The files should be compressed before transmission, and decompressed upon retrieval by the client, so that the transmission is as quick as possible, and as little space as possible is taken up on the server,
    \item Only text files need to be compressed, as most other files are already compressed,
    \item There should be no graphical user interface, but instead a command-line interface that can be used non-interactively,
    \item There is a hard deadline of the 31\textsuperscript{st} March 2026 for finishing the project.
\end{itemize}

\subsection{Specification}

Based on the requirements established in the client interview, I have established the following specification points. These will guide me throughout the development process, and can serve as a metric against which the final product is evaluated.

\begin{enumerate}
    \item The program should allow the user to create file backups on a server, run by the client themselves, 
    \begin{enumerate}
        \item The software should be able to run on both the client and server, with appropriate functionality for each situation,
        \item They should be able to work together, with a server being able to store the user’s backed up files, 
    \end{enumerate}
    \item The client program should run on modern Linux systems, usable on a laptop or desktop computer, 
    \begin{enumerate}
        \item An installable binary package should be created that is executable on amd64 Linux systems, 
        \item The client should not at any point use more than 500MB of RAM,
        \item The client should never use more than 5\% of CPU time while running on a modern laptop; my client's processor, the Intel i9-12950HX, can be used for a test case,
    \end{enumerate}
    \item It should be usable with an intuitive command-line interface, not requiring interactive input, 
    \begin{enumerate}
        \item All required functions of the program should be able to be passed in through command line flags and parameters, 
        \item There can be an optional interactive mode, but this should not allow any functionality than the non-interactive mode does not, 
        \item There should be no graphical interface, 
    \end{enumerate}
    \item The server should be able to run on a low-end laptop 10-15 years old, 
    \begin{enumerate}
        \item The program should avoid reading and writing to the system drives as much as possible, since most such devices use HDDs, 
        \item The server program should not at any point use more than 250MB of RAM
        \item The server program should never use more than 10\% of CPU time while running on such a device; my client's desired server's processor, a 2\textsuperscript{nd} generation Intel i5 processor, can be used as a test case, 
    \end{enumerate}
    \item All file transfers should be secure, so a man-in-the-middle would not be able to intercept data,
    \begin{enumerate}
        \item Any data intercepted should be of a format that the interceptor would not be able to decipher, or gain any useful information from, 
        \item The methods used should have their correctness checked,
    \end{enumerate}
    \item Text files stored on the server should be compressed such that they use a reasonable amount of storage, 
    \begin{enumerate}
        \item They should be at most as large as the uncompressed files, i.e. file size should never increase following compression, 
        \item The client should be able to decompress any file when needed by the user, storing any needed details about the compression within the file, 
    \end{enumerate}
    \item File transfers should be verifiably correct, so data is not corrupted during transmission,
    \begin{enumerate}
        \item Once a file has been transferred to the server, it should be compared with the original file on the client device, 
        \item This checking should be with a much smaller version of the file than that which was originally sent for most files – it should be up to 10kB in size, 
    \end{enumerate}
    \item Generally, the program should work robustly,
        \begin{enumerate}
            \item Problems with user input or command line arguments should be gracefully handled,
            \item The program should not crash while running,
            \item Errors should be reported to the client, server, or both, depending on which is appropriate in the situation.
        \end{enumerate}
\end{enumerate}

\subsection{Technical Points}

Based on the specification points established, I have extracted the key technical challenges that will be involved in building the software below. 

\begin{enumerate}
    \item Implement a transmission system with asymmetric encryption for key exchange between the client and the server, 
    \item Implement symmetric encryption between the client and the server after keys have been established, 
    \item Use a method to verify the integrity of the files transferred, 
    \item Implement a lossless compression algorithm for the text files stored on the server. 
\end{enumerate}

\subsection{Program Structure Diagram}

This diagram (Figure \ref{ProgramStructureDiagram}) shows how these main technical blocks will interact with each other in the client and server programs. While they will run different main procedures, these will share some key components.

\begin{figure}[ht]
    \centering
    \includegraphics[width=1.00\linewidth]{images/Analysis-program-structure.jpg} 
    \caption{Program Structure Diagram}
    \label{ProgramStructureDiagram}
\end{figure}

\clearpage

I will now consider how I can implement these main blocks.

\subsection{Transmission Procedure}

This is one of the largest components of the program, which is used for any and all communications between the client and server devices. There are some key points that need to be individually considered:
\begin{enumerate}
    \item  There needs to be a procedure used to exchange keys between the sender and receiver, perhaps in a similar style to the RSA algorithm, providing asymmetric encryption, 
    \item There must be a symmetric encryption algorithm used to encrypt communications following the establishment of keys, 
    \item The communications should be verified to be correct, using a method to check that there has not been any corruption during data transfer.
\end{enumerate}

\subsubsection{Key Exchange} \label{KeyExchangeSSS}

The three algorithms I will consider for this are RSA, Diffie-Hellman Key Exchange, and Elliptic Curve Cryptography.

These algorithms are commonly used, and very well documented. All three rely on mathematical problems that are intractable to solve without some prior knowledge: RSA relies on the difficulty in factoring a large semiprime integer (one with exactly two prime factors), while Diffie-Hellman and elliptic curve systems rely on the discrete logarithm problem. 

While elliptic curve cryptography can provide equivalent security with much smaller keys, and requires less processing power for the sender and receiver, this is unlikely to be at all impactful on the small scale of this backup utility, where this asymmetric encryption is just used to share a symmetric encryption key. However, these algorithms are tough to implement, and involve very high-level mathematics. This could increase the time taken to deliver the program past the client’s deadline, and adds more avenues for potential bugs to develop. This can be very dangerous, as shown in the following paper \cite{10.1007/978-3-642-27954-6_11}, where a small error in the implementation led to researchers being able to extract the full key. 

Plus, implementing it could cause me to fall short of my deadline, and it is very important for my client that I do not. Therefore, I will exclude elliptic curve cryptography from my choices.

I will now consider the basic implementation details of the RSA and Diffie-Hellman algorithms to aid my decision-making.

\bigskip
\textbf{RSA}
 
The website \cite{RSAgfg} contains details about the mathematics of the RSA algorithm, and how it can be implemented. From this, I have generated the following steps that are taken in the RSA algorithm.
\begin{enumerate}
    \item The receiver generates two large primes, $p,q$, then computes their product, $n = pq$.
    \item Euler’s Totient function, $\phi(pq) = (p-1)(q-1)$, is applied to $n$.
    \item An encryption exponent, $e$, is selected such that $1 < e < \phi(n)$, and $e$ is relatively prime with $\phi(n)$. Note that $a,b \in \mathbb{Z}^+$ are relatively prime $\iff \nexists\ x \in \mathbb{Z} : x > 1,\ x\,|\,a,$ and $x\,|\,b$.
    \item A decryption exponent, $d$, is calculated such that $de = 1 \pmod{\phi(n)}$. The extended Euclidian algorithm provides a method for this.
    \item The public key is formed from $(n,e)$ and the private key $(n,d)$. The receiver sends their public key only to the sender. 
    \item The encrypted ciphertext, $C$, is calculated through $C = M^e \pmod n$, where $M$ is some numerical representation of the message. 
    \item The ciphertext is decrypted into the plaintext message, $M$, by the calculation $M = C^d \pmod n$.
\end{enumerate}

Now, the sender encrypts their message with the receiver’s public key, and then only the receiver can decrypt it since only they have the private key. Computing $\phi(n)$ is infeasible without knowing the primes, $p,q$, that $n$ is a product of, and this is also infeasible to compute due to the large size of the primes. This means no other party can calculate the private key from the public key.

\bigskip
\textbf{Diffie-Hellman key exchange}

Page 84 of the following lecture notes \cite{OxCompSecNotes} contains a brief description of the mathematics behind the Diffie-Hellman key exchange. Below I describe the key steps in this algorithm in my own words.
\begin{enumerate}
    \item Either the sender or the receiver randomly generates a large prime, $p$, and a large number, $g$, and sends these to the other. 
    \item The sender and receiver then generate random secrets $s_a$ and $s_b$ respectively, for $s_a,s_b \in \mathbb{Z}^+$. 
    \item The sender sends $a = g^{s_a}\pmod p$ to the receiver, who computes $a^{s_b} = g^{s_as_b}\pmod p$. 
    \item The receiver sends $b = g^{s_b}\pmod p$ to the receiver, who computes $b^{s_a} = g^{s_as_b}\pmod p$. 
    \item Now, they both have the shared secret, $S = g^{s_as_b}\pmod p$, which can be used to establish a shared key.
\end{enumerate}

An attacker would have access to $g,\ p,\ g^{s_a}\pmod p,\ g^{s_b}\pmod p$, and need to compute $S = g^{s_as_b}\pmod p$ to compromise the conversation. A quick solution to the Discrete Logarithm Problem would imply a quick solution to this problem, but one is not currently known. This is central to the security of this key exchange, similarly to the prime factorisation problem for RSA. 

As can be seen, RSA and Diffie-Hellman are very similar in security, and difficulty of implementation, however the advantage of RSA is that it can also be used to provide signed certificates to prove the authenticity of the sender. This has not been an aim as of yet, but using RSA will make it easier to implement this if desired in the long-term, since I intend to make this project open-source upon completion.

\subsubsection{Symmetric Encryption}

Following the establishment of keys, it is greatly advantageous to use a symmetric encryption system rather than continuing with asymmetric systems, since these are much faster and more efficient. A similar process is used by the TLS protocol, which secures more internet communications. 

There are a great number of available symmetric encryption algorithms, including the Caesar Cipher, AES, ChaCha20, and Blowfish, among others.

The Caesar cipher can be immediately eliminated, since it is trivial to break with even very simple cryptanalysis. However, the other algorithms are very similar in that they all provide secure encryption that is very tough to break. The final choice of algorithm should be one that is well-documented, simple to implement, and efficient to run.

However, AES has been used as the basis of internet security for decades, and as such there has been great research into verifying its security, and ensuring implementations do not contain vulnerabilities. Plus, C\# provides a standard implementation of this. Therefore, I will use AES for my symmetric encryption needs.

\subsubsection{Verifying the Integrity of Communications} \label{VerifyingIntegritySSS}

It is necessary to ensure all communications have arrived at the desired recipient uncorrupted, exactly as they were sent. Therefore, some method must be used to ensure that this is the case. 

Some simple methods for this include majority voting, where each bit is sent an odd number ($> 1$) of times, and the most common bit is accepted, and parity bits, where a bit is added to the end of each string of defined length to ensure that the number of 0s or 1s is either even or odd as desired. 

However, these methods allow multiple mistakes to “cancel” each other out, and say little about the correctness of the whole message. A more robust method is to compute some form of hash from the message, and send that directly after it. The receiver will also compute a hash of the data received, and compare it to the hash sent. If these match, then the data has arrived intact. This verifies that the whole message has arrived safely, while requiring comparatively little additional data to be transferred. Plus, these algorithms can run very quickly.

I can use the SHA-256 algorithm for this. It is a cryptographically secure hash function, that produces a 256-bit hash for any length of text. It is widely used, and it is very unlikely for two distinct inputs to give the same output. There is also an available C\# implementation that I can use for this.

\subsection{Compression} \label{CompressionAnalysis}

Firstly, it is important to establish that I will want to use a lossless algorithm, since a backup program should not change the files being backed up – the whole point is to provide a safe copy of the actual system files themselves. Plus, lossy compression of important system files would be catastrophic. 

There is also the key problem that many files being backed up will already be compressed. The system must be able to deal with these, not producing output with larger file size than the uncompressed versions. 

There is also the important consideration of whether each file is individually compressed, or whether the whole block of files is compressed together. 

There are two broad options I can take based on the above. I could either compress the whole block of data, or each file individually, with some metadata at the beginning of the file outlining what procedure (if any) was used.

The first option will not be very useful, since different file types can be effectively compressed using vastly different algorithms, so I will not be able to achieve significant size reductions in this way. Therefore, I believe the right choice is to compress files individually. 

This means that I would need to write different algorithms for different file types, and have a way of detecting files that have already been compressed. However, the discussion with my client revealed that he already compresses the majority of non-text files on his system. So, I have decided to just implement compression for plaintext files, where the most space savings can be made, while minimising the unnecessary processing taking place, so the program executes as quickly as possible. This will also reduce development time, allowing me to finish the program in time for my client’s strict deadline. 

Three lossless plaintext compression algorithms I will consider as the basis of my algorithm are Run-Length Encoding, Huffman Coding, and Prediction by Partial Matching.

Run-Length Encoding works by representing repeated runs of one character as frequency-data pairs. For example, AAAAAABBB would be stored as "6A3B". However, this is not very effective for text files, where there is not often repeated runs of characters; it is instead very useful for images, where there may be, for example, a dark background with many long runs of dark pixels. Therefore, it is unsuitable for my needs, to compress plaintext.

Prediction by Partial Matching uses complex statistical techniques to predict the most probable next characters based on the previous ones, assigning more probable characters shorter binary representations. However, to create and optimise such an algorithm, it would take a long time, which may cause be to pass my client's non-negotiable deadline. Plus, it is more computationally complex than the alternatives, so the compression and decompression algorithms would use more memory and processor time, and would likely take longer to run.

Huffman Coding provides a perfect balance between complexity and how effective it is at compressing plaintext. It looks at the file as a whole, assigning more frequently used characters shorter binary representations. This is very effective for text files, where some letters are used more often than others. Moreover, most available Unicode characters are never used (special characters, scientific symbols, foreign language characters...) in most files, so the majority of characters would have a shorter representation in a file compressed in this way. While it may be more complex than Run-Length Encoding, it is perfectly feasible for me to implement it in the given time, and the implementations can be very quick! Even though Prediction by Partial Matching can in some cases provide better compression, for the files where this really matters to my client, they can just compress them themselves before backing them up.

I have researched the Huffman Coding algorithm through this website \cite{PurdueHuffman}, which gives an explanation of one implementation of it.

The high-level steps of the encoding algorithm are as follows:
\begin{enumerate}
    \item Count the number of times each character appears in the text,
    \item Turn each character into a tree with a single node, weighted by its frequency,
    \item Continually merge the two nodes with the smallest weights into a new tree, creating a new node with the nodes being merged as children, with its weighting being the sum of its two children's weights,
    \item Stop once there is only one tree left - this is the Huffman tree,
    \item Carry out a post-order tree traversal on the Huffman tree to write a representation of it to the beginning of the compressed file, and to build up a mapping of characters to their binary Huffman representation (this is defined by their placement in the tree with respect to the root node),
    \item Go through each character in the file, writing its Huffman representation to the compressed file.
\end{enumerate}

The high level steps of the decoding algorithm are as follows:
\begin{enumerate}
    \item Read the Huffman tree from the start of the compressed file, representing it as a tree in the program,
    \item Carry out a post-order traversal of the tree to build up a mapping of the Huffman binary bit-strings to the characters they represent,
    \item Go through all of the bits in the compressed file, writing the character they represent to the decompressed file.
\end{enumerate}

\subsection{Class Diagram}

Figure \ref{AnalysisClassDiagram} shows a proposed design for the classes in my program.

\begin{figure}[ht]
    \centering
    \includegraphics[width=\linewidth]{images/Analysis-class-diagram.jpg} 
    \caption{Proposed UML Class Diagram}
    \label{AnalysisClassDiagram}
\end{figure}

The FileSender and FileReceiver classes will contain methods related to sending and receiving data. These methods can be used by both the client and server programs, and so the Client and Server classes - which will contain the main methods run by the client and the server respectively - depend on them. These classes both inherit from an abstract base FileTransferer class which will contain any methods or attributes common to both classes.

The HuffmanNode is a custom data structure representing the nodes that form a Huffman tree. This is used by both the HuffmanEncoder and HuffmanDecoder classes, which contain methods for encoding plaintext into Huffman, and decoding Huffman compressed text back into plaintext respectively. The ByteMaker and ByteReader classes will deal with any complexities that may arise in the processor of reading from and writing to binary files.

Finally, the SetupRSA class will contain methods and data relating to the creation of RSA public and private keys, SetupAES will establish the shared AES encryption key, the Hasing class will allow methods for generating and verifying hashes, and the Transmission class will contain common procedures needed to serialise data into a format that can be decoded once received, and to decode such data once it has been received. These are of use in the FileSender and FileReceiver classes.

\subsection{Order of Completion}

This is a rough timeline of how long I expect it would take me to complete the sections outlined above.
\begin{enumerate}
    \item Plaintext compression - 2 weeks
    \item RSA style key exchange - 2 weeks 
    \item Hashing to verify transmission, and symmetric encryption to secure it - 1 week 
    \item Finalising the transmission system - 2 weeks
    \item Terminal interface - 1 week 
    \item Putting it all together - 2 weeks
\end{enumerate}

In total, this will take 10 weeks, leaving me on track to meet my deadline with plenty of time to spare in case unexpected issues arise.

\end{document}
